\documentclass[11pt]{article}

% Change "review" to "final" to generate the final (sometimes called camera-ready) version.
% Change to "preprint" to generate a non-anonymous version with page numbers.
\usepackage[]{acl}
\usepackage{times}
\usepackage{latexsym}

% For proper rendering and hyphenation of words containing Latin characters (including in bib files)
\usepackage[T1]{fontenc}
% For Vietnamese characters
% \usepackage[T5]{fontenc}
% See https://www.latex-project.org/help/documentation/encguide.pdf for other character sets

% This assumes your files are encoded as UTF8
\usepackage[utf8]{inputenc}
\usepackage{microtype}
\usepackage{inconsolata}
\usepackage{graphicx}

% If the title and author information does not fit in the area allocated, uncomment the following
%
%\setlength\titlebox{<dim>}
%
% and set <dim> to something 5cm or larger.

\title{Group X Progress Report:\\My Group's Project Name}


\author{Bhavya Patel, Second Author, Third Author \\
  \texttt{\{patelb61,macid2,macid3\}@mcmaster.ca} }

%\author{
%  \textbf{First Author\textsuperscript{1}},
%  \textbf{Second Author\textsuperscript{1,2}},
%  \textbf{Third T. Author\textsuperscript{1}},
%  \textbf{Fourth Author\textsuperscript{1}},
%\\
%  \textbf{Fifth Author\textsuperscript{1,2}},
%  \textbf{Sixth Author\textsuperscript{1}},
%  \textbf{Seventh Author\textsuperscript{1}},
%  \textbf{Eighth Author \textsuperscript{1,2,3,4}},
%\\
%  \textbf{Ninth Author\textsuperscript{1}},
%  \textbf{Tenth Author\textsuperscript{1}},
%  \textbf{Eleventh E. Author\textsuperscript{1,2,3,4,5}},
%  \textbf{Twelfth Author\textsuperscript{1}},
%\\
%  \textbf{Thirteenth Author\textsuperscript{3}},
%  \textbf{Fourteenth F. Author\textsuperscript{2,4}},
%  \textbf{Fifteenth Author\textsuperscript{1}},
%  \textbf{Sixteenth Author\textsuperscript{1}},
%\\
%  \textbf{Seventeenth S. Author\textsuperscript{4,5}},
%  \textbf{Eighteenth Author\textsuperscript{3,4}},
%  \textbf{Nineteenth N. Author\textsuperscript{2,5}},
%  \textbf{Twentieth Author\textsuperscript{1}}
%\\
%\\
%  \textsuperscript{1}Affiliation 1,
%  \textsuperscript{2}Affiliation 2,
%  \textsuperscript{3}Affiliation 3,
%  \textsuperscript{4}Affiliation 4,
%  \textsuperscript{5}Affiliation 5
%\\
%  \small{
%    \textbf{Correspondence:} \href{mailto:email@domain}{email@domain}
%  }
%}

\begin{document}
\maketitle
% \begin{abstract}
% \end{abstract}

\section{Introduction}

Here, write a brief introduction to the problem you are solving. This can be adapted from your problem description and motivation from the original proposal. This should be around 0.25-0.5 pages.

\section{Related Work}

Here, talk about the related work you encountered for your approach. Cite at least 5 references. Refer to item 2. No one has done exactly your task? Write about the most similar thing you can find. This should be around 0.25-0.5 pages.

\section{Dataset}

We are using the publicly available Myocardial Infarction Complications dataset from the UCI Machine Learning Repository. As the dataset is annotated, no additional annotation was done. The dataset contains 1700 patient records with 111 features and 12 target columns. The input features include patient demographics along with clinical measurements and information. The 12 target variables consist of 11 binary outcomes and one 8-class categorical outcome.

\subsection{Feature Selection and Missing Data}

During the initial data analysis, columns corresponding to information available after 48 hours of hospitalization (\texttt{R\_AB\_3\_n}, \texttt{NA\_R\_3\_n}, and \texttt{NOT\_NA\_3\_n}) were dropped, as it did not align with our goal. The missing data distributions for each feature were analyzed and anything with over 50\% missing data was examined closely. Features \texttt{KFK\_BLOOD} and \texttt{IBS\_NASL} were dropped immediately as they had 99.76\% and 95.76\% missing values, respectively.

Two features representing emergency response blood pressure measurements (\texttt{D\_AD\_KBRIG} and \texttt{S\_AD\_KBRIG}) had 63.29\% missing values. Correlation analysis was run between these columns and their corresponding ICU blood pressure measurements: the correlation between \texttt{D\_AD\_KBRIG} and \texttt{D\_AD\_ORIT} was 0.56 while the correlation between \texttt{S\_AD\_KBRIG} vs \texttt{S\_AD\_ORIT} was 0.61. They were not strongly correlated enough to justify removal. A missing analysis was conducted and both columns had missing values for the same instances, so one binary indicator column was created (\texttt{EMR\_RES\_BLOOD}) to represent both columns. \texttt{EMR\_RES\_BLOOD} would have a value of 1 if ER blood pressure measurement was taken or 0 if not taken.

After feature selection was complete, there were 107 input features remaining. These features were analyzed and divided into these categories: numerical, binary, ordinal or partially ordered to determine how to appropriately preprocess them. A correlation analysis was run for the numerical features, however most of the relationships ranged from weak to moderate, and therefore no further features were dropped.

\subsection{Target Distribution and Data Splitting}

Although the target columns contained no missing values, there was a large class imbalance. Most binary target columns had above 1600 ``no'' instances and under 100 ``yes'' instances. The categorical target column ``LET\_IS'' had categories 6 and 7 representing under 2\% of the outcomes. To preserve the distribution of possible outcomes across the train, test and validation sets, the split was done using iterative stratification. As iterative stratification required binary labels, the categorical target column LET\_IS was temporarily one-hot encoded into 8 columns. The indices found in the split were then applied to do the split on the original dataset to allow LET\_IS to remain a categorical column. The dataset was divided into 70\% train, 15\% validation and 15\% test.

\subsection{Missing Value Imputation and Feature Scaling}

Missing value imputation and feature scaling were done after the split to prevent data leakage. The numerical features were imputed using train medians as they are less sensitive to outliers. The categorical columns (binary, ordinal and partially-ordered) columns were imputed using the train modes. These imputation values were then applied to test and validation sets.

For feature scaling, similarly, the scalers were fitted to the train set before being applied to the validation and test sets. The numerical features were scaled using \texttt{StandardScaler}. The ordinal features and the partially ordered feature were scaled using \texttt{MinMaxScaler} to map values between 0 and 1. The binary columns were left as is. The resulting datasets were named: \texttt{X\_train}, \texttt{y\_train}, \texttt{X\_validation}, \texttt{y\_validation}, \texttt{X\_test}, \texttt{y\_test}.

\subsection{Future Considerations}

In the future, the partially ordered column \texttt{ZSN\_A} might be investigated further as it has 2 severity branches and therefore does not follow a single ordinal ranking. It might be split into 2 features, each representing a ``severity branch''. For the current preprocessing, it has been left in its original form.

\section{Features}

Describe any features you used for your model, or how your data was input to your model. Are you doing feature engineering or feature selection? Are you learning embeddings? Is it all part of one neural network? Refer to item 2. This may range from 0.25 pages to 0.5 pages.

\section{Implementation}

Describe your model and implementation here. Refer to item 4. This may take around a page.

\section{Results and Evaluation}

How are you evaluating your model? What results do you have so far? What are your baselines? Refer to item 5. This may take around 0.5 pages.

\section{Feedback and Plans}

The TA (Shreya) recommended comparing our current MLP using 48-hour features with the same architecture using 72-hour features to see the impact of reduced input features. She also suggested experimenting with L1 and L2 regularization to reduce overfitting, exploring other models, and including all project files with a detailed README.

Our initial experiments addressed class imbalance using full inverse-frequency, square-root weighting, and weighting with an exponent of 0.7. Full weighting produced many false positives (127 for \texttt{PREDS\_TAH} and 0 true positives), while square-root weighting predicted almost all negatives. The 0.7 weighting performed better than the baseline on some targets, but all models struggled with rare complications. The full weighting model also showed a lower training loss (under 2) than validation loss (around 3.3), suggesting possible overfitting.

Moving forward, we plan to evaluate binary targets using precision, recall, F1-score, and precision-recall AUC as well as doing performance evaluation for each class in the \texttt{LET\_IS} target column. We will investigate L1/L2 regularization, alternative architectures, and cross-validation due to the low number of positive outcomes. After selecting the best configuration using validation results, we will evaluate it on the test set. Finally, we will compare the 48-hour and 72-hour feature sets to determine whether complications can be predicted using only information available within the first 48 hours of hospitalization which addresses our main research question.

\iffalse
\section{Template Notes}

You can remove this section or comment it out, as it only contains instructions for how to use this template. You may use subsections in your document as you find appropriate.

\subsection{Tables and figures}

See Table~\ref{citation-guide} for an example of a table and its caption.
See Figure~\ref{fig:experiments} for an example of a figure and its caption.


\begin{figure}[t]
  \includegraphics[width=\columnwidth]{example-image-golden}
  \caption{A figure with a caption that runs for more than one line.
    Example image is usually available through the \texttt{mwe} package
    without even mentioning it in the preamble.}
  \label{fig:experiments}
\end{figure}

\begin{figure*}[t]
  \includegraphics[width=0.48\linewidth]{example-image-a} \hfill
  \includegraphics[width=0.48\linewidth]{example-image-b}
  \caption {A minimal working example to demonstrate how to place
    two images side-by-side.}
\end{figure*}


\subsection{Citations}

\begin{table*}
  \centering
  \begin{tabular}{lll}
    \hline
    \textbf{Output}           & \textbf{natbib command} & \textbf{ACL only command} \\
    \hline
    \citep{Gusfield:97}       & \verb|\citep|           &                           \\
    \citealp{Gusfield:97}     & \verb|\citealp|         &                           \\
    \citet{Gusfield:97}       & \verb|\citet|           &                           \\
    \citeyearpar{Gusfield:97} & \verb|\citeyearpar|     &                           \\
    \citeposs{Gusfield:97}    &                         & \verb|\citeposs|          \\
    \hline
  \end{tabular}
  \caption{\label{citation-guide}
    Citation commands supported by the style file.
  }
\end{table*}

Table~\ref{citation-guide} shows the syntax supported by the style files.
We encourage you to use the natbib styles.
You can use the command \verb|\citet| (cite in text) to get ``author (year)'' citations, like this citation to a paper by \citet{Gusfield:97}.
You can use the command \verb|\citep| (cite in parentheses) to get ``(author, year)'' citations \citep{Gusfield:97}.
You can use the command \verb|\citealp| (alternative cite without parentheses) to get ``author, year'' citations, which is useful for using citations within parentheses (e.g. \citealp{Gusfield:97}).

\subsection{References}

\nocite{Ando2005,andrew2007scalable,rasooli-tetrault-2015}

Many websites where you can find academic papers also allow you to export a bib file for citation or bib formatted entry. Copy this into the \texttt{custom.bib} and you will be able to cite the paper in the \LaTeX{}. You can remove the example entries.

\subsection{Equations}

An example equation is shown below:
\begin{equation}
  \label{eq:example}
  A = \pi r^2
\end{equation}

Labels for equation numbers, sections, subsections, figures and tables
are all defined with the \verb|\label{label}| command and cross references
to them are made with the \verb|\ref{label}| command.
This an example cross-reference to Equation~\ref{eq:example}. You can also write equations inline, like this: $A=\pi r^2$.


% \section*{Limitations}
\fi

\section*{Team Contributions}

Bhavya Patel worked on intial data analysis and the data preprocessing. She wrote sections: "Dataset" and "Feedback and Plans".
% Bibliography entries for the entire Anthology, followed by custom entries
%\bibliography{custom,anthology-overleaf-1,anthology-overleaf-2}

% Custom bibliography entries only
\bibliography{custom}

% \appendix

% \section{Example Appendix}
% \label{sec:appendix}

% This is an appendix.

\end{document}
